\section*{NeurIPS Paper Checklist}

\begin{enumerate}

\item {\bf Claims}
    \item[] Question: Do the main claims made in the abstract and introduction accurately reflect the paper's contributions and scope?
    \item[] Answer: \answerYes{}
    \item[] Justification: The main claims are formulated explicitly in Section 1 (Core Research Questions RQ1--RQ4, Summary of Contributions) and supported by empirical benchmarks in Section 4 and formal proofs in Section 2 and Appendix A.

\item {\bf Limitations}
    \item[] Question: Does the paper discuss the limitations of the work performed by the authors?
    \item[] Answer: \answerYes{}
    \item[] Justification: Section 2, Section 3, and Appendix A disclose all architectural parameters, 3D BFS masking parameters, optimization schedules, loss weights, and deterministic seeding protocols needed for exact reproduction. The complete codebase is packaged in \texttt{thesis\_3d}.
    \item[] Guidelines:
    \begin{itemize}
        \item The answer \answerNA{} means that the paper does not include experiments.
        \item If the paper includes experiments, a \answerNo{} answer to this question will not be perceived well by the reviewers: Making the paper reproducible is important, regardless of whether the code and data are provided or not.
        \item If the contribution is a dataset and\slash or model, the authors should describe the steps taken to make their results reproducible or verifiable. 
        \item Depending on the contribution, reproducibility can be accomplished in various ways. For example, if the contribution is a novel architecture, describing the architecture fully might suffice, or if the contribution is a specific model and empirical evaluation, it may be necessary to either make it possible for others to replicate the model with the same dataset, or provide access to the model. In general. releasing code and data is often one good way to accomplish this, but reproducibility can also be provided via detailed instructions for how to replicate the results, access to a hosted model (e.g., in the case of a large language model), releasing of a model checkpoint, or other means that are appropriate to the research performed.
        \item While NeurIPS does not require releasing code, the conference does require all submissions to provide some reasonable avenue for reproducibility, which may depend on the nature of the contribution. For example
        \begin{enumerate}
            \item The code is made available as part of the submission or is promised to be released upon acceptance.
            \item The data is made available as part of the submission or is promised to be released upon acceptance.
            \item A pipeline for recreating the data is made available as part of the submission or is promised to be released upon acceptance.
        \end{enumerate}
    \end{itemize}

\item {\bf Open access to data and code}
    \item[] Question: Does the paper provide open access to data and code?
    \item[] Answer: \answerYes{}
    \item[] Justification: The code repository \texttt{thesis\_3d} is fully standalone with automated pytest suites, reproducible configs, CLI runners, and Kaggle runner notebooks. The BraTS 2024 GLI challenge dataset is publicly available through Synapse.

\item {\bf Experimental setting/details}
    \item[] Question: Does the paper specify all the training and evaluation details?
    \item[] Answer: \answerYes{}
    \item[] Justification: Section 3 details optimizer choices (AdamW), learning rates, warmup schedules, weight decays, batch sizes, mixed-precision settings, and evaluation metrics.

\item {\bf Experiment statistical significance}
    \item[] Question: Does the paper report error bars or the statistical significance of the results?
    \item[] Answer: \answerYes{}
    \item[] Justification: All volumetric metrics (3D Dice, IoU, cKDTree HD95) are reported with sample means and standard deviations across the held-out test cohort.

\item {\bf Experiments compute resources}
    \item[] Question: For each experiment, does the paper provide sufficient information on the computer resources needed to reproduce the experiments?
    \item[] Answer: \answerYes{}
    \item[] Justification: Section 3.2 and Section 4.1 report hardware specifications (NVIDIA Tesla T4 16GB / Apple Silicon), training durations per epoch, peak VRAM allocations, and forward inference latencies (in ms per $128^3$ volume).

\item {\bf Code of ethics}
    \item[] Question: Does the research conform with the NeurIPS Code of Ethics?
    \item[] Answer: \answerYes{}
    \item[] Justification: The research uses openly published, fully de-identified clinical MRI challenge data (BraTS 2024 GLI) and strictly conforms with institutional ethics and the NeurIPS Code of Ethics.

\item {\bf Broader impacts}
    \item[] Question: Does the paper discuss potential positive and negative societal impacts?
    \item[] Answer: \answerYes{}
    \item[] Justification: Section 5.3 outlines the clinical utility of label-efficient 3D segmentation in resource-constrained radiotherapy planning, while addressing clinical safety boundaries against unsupervised off-label diagnostic deployment.

\item {\bf Safeguards}
    \item[] Question: Does the paper describe safeguards for responsible release of models?
    \item[] Answer: \answerNA{}
    \item[] Justification: The models are discriminative volumetric encoders dedicated to brain glioma MRI segmentation and pose no dual-use or disinformation risks.

\item {\bf Licenses for existing assets}
    \item[] Question: Are creators of existing assets properly credited and licenses respected?
    \item[] Answer: \answerYes{}
    \item[] Justification: Open benchmarks (BraTS 2024 GLI) and open-source libraries (PyTorch, MONAI, NiBabel) are cited with proper attributions.

\item {\bf New assets}
    \item[] Question: Are new assets introduced in the paper well documented?
    \item[] Answer: \answerNA{}
    \item[] Justification: The paper utilizes existing public benchmark datasets.

\item {\bf Crowdsourcing and research with human subjects}
    \item[] Question: Does the paper involve crowdsourcing or research with human subjects?
    \item[] Answer: \answerNA{}
    \item[] Justification: No crowdsourcing or human subject experiments were performed.

\item {\bf Institutional review board (IRB) approvals}
    \item[] Question: Does the paper describe IRB approvals for research with human subjects?
    \item[] Answer: \answerNA{}
    \item[] Justification: The research exclusively utilizes previously published, de-identified public challenge datasets with pre-existing ethics approvals.

\item {\bf Declaration of LLM usage}
    \item[] Question: Does the paper describe the usage of LLMs?
    \item[] Answer: \answerNA{}
    \item[] Justification: Large language models were not used as a methodological component of the algorithmic design or experimental pipeline.

\end{enumerate}
